% SLICESETTRANS1 2026-10-16 09:15 Asia/Seoul
\documentclass[11pt]{article}
\usepackage[T1]{fontenc}
\usepackage{lmodern,amsmath,amssymb,amsthm,booktabs,hyperref}
\usepackage[margin=0.9in]{geometry}
\newtheorem{theorem}{Theorem}
\newtheorem{lemma}[theorem]{Lemma}
\newtheorem{corollary}[theorem]{Corollary}
\newcommand{\EXACT}{\textbf{PROVED/EXACT}}
\newcommand{\KNOWN}{\textbf{KNOWN THEOREM}}
\newcommand{\COMP}{\textbf{COMPUTATIONAL}}
\newcommand{\OPEN}{\textbf{OPEN}}
\newcommand{\RET}{\textbf{FALSIFIED/RETIRED}}
\title{SLICESETTRANS1: Set-Circuit and Transversal Formulation of the Exact-Density UNSAT Slice}
\author{P versus NP project}
\date{16 October 2026}
\begin{document}\maketitle

\noindent\textbf{Status.} The reduction of the exact-density UNSAT slice to a Wegener set-circuit target, the transversal identities, the Hamming tensor-rank theorem, the additive-potential ceiling, and the masked-KW polynomial DAG protocol are \EXACT at their stated models, pending independent review. Wegener's set-circuit equivalence, Hrube\v{s}--Pudl\'ak's full-cube monotone lower bound, Hiltgen--Paterson mass production, Dunne's central-slice results, and Sokolov's DAG-like KW characterization are \KNOWN. The superpolynomial set-circuit lower bound for the exact UNSAT transversal family remains \OPEN; hence $P\ne NP$ remains \OPEN.

\section{IDEAS PROPOSED BEFORE LITERATURE}
\paragraph{A: native set-circuit target.} Replace the distributional approximation language by Wegener's exact set-circuit representation of a slice. Prediction: the output family is exactly the family of exact-$3q$ UNSAT formulas, equivalently exact-cardinality transversals of assignment-falsification blocks.
\paragraph{B: transfer full UNSAT monotone hardness.} Hrube\v{s}--Pudl\'ak note exponential monotone lower bounds for full $\mathrm{UNSAT}_{3,q}$. Test whether Dunne/Wegener central-slice reductions preserve that monotone hardness. Falsifier: the known slice-power barrier or the constant-weight monotone-map theorem.
\paragraph{C: HD/Hamming association space.} Use the assignment-by-clause incidence matrix and its tensor powers as an exact spectral coordinate system. Missing theorem: a gate-stable potential for unions/intersections.
\paragraph{D: DAG-like KW route.} Translate set-circuit size into DAG-like communication size and search for a hard restricted relation. Falsifier: a polynomial DAG protocol for the natural assignment-witness restriction.

\section{LITERATURE CHECK AND REVISIONS}
Wegener's set-circuit theorem states that for a $k$-slice $f$, if $SC(f)$ is the minimum number of union/intersection gates needed to compute the set of length-$k$ prime implicants from the coordinate stars, then $SC(f)\le C_m(f)$ and conversely monotone complexity is at most $SC(f)$ plus threshold/pseudo-input overhead. Thus a superpolynomial set-circuit lower bound is sufficient for the desired monotone slice lower bound.

Hrube\v{s}--Pudl\'ak explicitly state that exponential monotone circuit lower bounds for full $\mathrm{UNSAT}_{3,q}$ follow from known monotone lower bounds. The same paper explicitly warns that their approximate criterion cannot restrict the soundness condition to one exact Hamming layer because monotone circuits for slice functions are as powerful as general circuits. Thus full-cube hardness is not silently inherited by our slice.

Dunne proves NP-completeness and circuit-equivalence phenomena for central slices of several NP-complete predicates, including a monotone encoding of SAT. His thesis states that every slice of those predicates is a projection of the central slice of a larger instance. This is useful complexity transport, but it is not a monotone projection from full clause-presence $\mathrm{UNSAT}_{3,q}$ into our fixed-weight proper-3CNF slice; the project constant-weight monotone-map theorem already forbids any nonconstant coordinatewise monotone full-cube encoding into a fixed layer.

Hiltgen--Paterson show that PI$_k$ set circuits admit strong mass-production sharing and use it to collapse a previously conjectured $\Theta(n^{3/2})$ Ne\v{c}iporuk-slice bound to linear size. Any proposed lower-bound invariant must therefore survive arbitrary DAG sharing rather than count minterms, local certificates, or independently generated outputs.

Sokolov proves a DAG-like analogue of Karchmer--Wigderson characterizing Boolean circuit size. This makes the full slice KW relation a valid alternate target, but the assignment-witness subrelation below is polynomially easy in DAG size and is retired.

\section{CONSTRUCTION OR PROOF ATTEMPT}
\subsection{Exact set-circuit target}
Let $N=8\binom q3$, $t=3q$, and let $U_{q,t}$ be zero below weight $t$, one above weight $t$, and equal to UNSAT on the exact layer. Put
\[\Omega_{q,t}=\binom{\mathcal C_q}{t},\qquad \mathcal U_{q,t}=\{F\in\Omega_{q,t}:F\text{ is UNSAT}\}.\]

\begin{theorem}[Prime-implicant identification]
The length-$t$ prime implicants of $U_{q,t}$ are exactly $\mathcal U_{q,t}$.
\end{theorem}
\begin{proof}
Every positive exact-$t$ input is minimal positive because every strict subset has weight below $t$ and is therefore zero. Conversely no monomial of length below $t$ can be an implicant because every input containing it can be completed to an input of weight below $t$ where the slice function is zero.
\end{proof}

For a clause coordinate $C$, Wegener's pseudo-input set is exactly the star
\[Y_C=\{F\in\Omega_{q,t}:C\in F\}.\]
Hence the load-bearing theorem can be stated without any hidden circuit representation:
\[\boxed{SC(\mathcal U_{q,t};\{Y_C\}_{C\in\mathcal C_q})=N^{\omega(1)}.}\]
By Wegener and Valiant, this one theorem would imply a superpolynomial unrestricted circuit lower bound for an explicit coNP-complete slice and therefore $P\ne NP$.

\subsection{Exact transversal geometry}
For an assignment $a\in\{0,1\}^q$, define
\[B_a=\{C\in\mathcal C_q:a\text{ falsifies }C\}.\]
Then $|B_a|=\binom q3$ and
\[F\in\mathcal U_{q,t}\iff \forall a,\ F\cap B_a\ne\varnothing.\]
Thus $\mathcal U_{q,t}$ is precisely the family of size-$t$ transversals of the hypergraph with vertex set $\mathcal C_q$ and hyperedges $B_a$.

\begin{lemma}[Hamming block intersections]
If $d=d_H(a,b)$, then
\[|B_a\cap B_b|=\binom{q-d}{3}.\]
More generally, if $a_1,\ldots,a_r$ agree simultaneously on exactly $e$ coordinates, then
\[\left|\bigcap_{j=1}^r B_{a_j}\right|=\binom e3.\]
\end{lemma}
\begin{proof}
A proper signed clause is falsified by all listed assignments exactly when its three variables are chosen among coordinates on which all assignments agree and its three signs are the unique signs falsified by that common restriction.
\end{proof}

The trivial exact set-circuit expression is
\[\mathcal U_{q,t}=\bigcap_{a\in\{0,1\}^q}\ \bigcup_{C\in B_a}Y_C,\]
an exponential upper bound. The lower-bound problem is to show that no polynomial DAG of unions/intersections can mass-produce the same transversal family.

\subsection{Exact HD/Hamming tensor space}
Let $A$ be the $2^q\times N$ incidence matrix
\[A_{a,C}=1[C\in B_a].\]
Each column is the indicator of one prescribed assignment on a 3-variable support.

\begin{theorem}[Assignment-block rank]
Over $\mathbb R$,
\[\operatorname{rank}A=\sum_{i=0}^3\binom qi.\]
More generally, if $\phi_r(a)=A_a^{\otimes r}$, then
\[\boxed{\operatorname{rank}\{\phi_r(a):a\in\{0,1\}^q\}=\sum_{i=0}^{\min\{3r,q\}}\binom qi.}\]
\end{theorem}
\begin{proof}
A clause-cylinder indicator has Fourier degree at most three, and the collection of all signed clauses spans every parity character of degree at most three. Products of $r$ such functions have degree at most $3r$. Conversely every parity character of degree at most $3r$ factors into at most $r$ characters of degree at most three, each in the span of the clause cylinders. Hence the tensor-coordinate span is exactly the Fourier subspace of degrees at most $3r$, whose dimension is the stated binomial sum.
\end{proof}

Consequently constant-order HD features have polynomial rank; $r=\Theta(\log q)$ gives dimension $\exp(\Theta((\log q)^2))$ up to lower-order terms; and $r\ge q/3$ gives the full $2^q$ assignment dimension. This is an exact asymptotic HD space, not a circuit lower bound.

\subsection{Why simple additive HD potentials cannot finish the proof}
Let $\Omega=\Omega_{q,t}$ and assign nonnegative weights $w_F$ to its atoms. For a family $S\subseteq\Omega$, put $W(S)=\sum_{F\in S}w_F$. Then
\[\sum_{C\in\mathcal C_q}W(Y_C)=tW(\Omega).\]
If the weights are supported on the target family, $W(\Omega)=W(\mathcal U_{q,t})$, so some input star already satisfies
\[W(Y_C)\ge \frac tN W(\mathcal U_{q,t}).\]
Since $W(S\cup T)\le W(S)+W(T)$ and $W(S\cap T)\le\min\{W(S),W(T)\}$, any lower-bound argument that only tracks multiplicative growth of this one nonnegative additive magnitude can extract at most
\[\log_2(N/t)=O(\log q)\]
gates from the output/input magnitude ratio. Therefore simple additive HD mass is \RET as a route to a superpolynomial set-circuit lower bound. Signed sketches do not repair gate stability: cancellations can make both parent sketches zero while the union sketch is nonzero.

\subsection{DAG-like KW target and an easy restricted relation}
The full monotone KW relation of the exact slice is:
\begin{quote}
Alice receives $F\in\mathcal U_{q,t}$, Bob receives satisfiable $G\in\Omega_{q,t}$, and they must output a clause $C\in F\setminus G$.
\end{quote}
By the DAG-like KW theorem, a superpolynomial lower bound on the appropriate rectangle-DAG size would lower-bound the monotone/set-circuit size.

A tempting restriction gives Bob an assignment $a$ and the set
\[H=F\cap S_a,\qquad S_a=\mathcal C_q\setminus B_a,\]
of all clauses of $F$ satisfied by $a$, with the task of outputting an element of $F\setminus H=F\cap B_a$. This restriction reduces locally to the full slice KW relation: Bob pads $H$ by canonical clauses from $S_a\setminus H$ until it has size $t$, obtaining an exact-$t$ formula $G$ satisfied by $a$; because the filler is outside $F$, $F\setminus G=F\setminus H$.

\begin{theorem}[Masked violation search has a polynomial DAG protocol]
On universe size $N$, with promise $H\subsetneq F$ and $|F|=t$, the search problem of finding an element of $F\setminus H$ has a deterministic binary DAG communication protocol of size $O(Nt\log t)$.
\end{theorem}
\begin{proof}
Fix a balanced binary partition tree of the coordinate universe. At an interval $I$, Alice communicates the count of $F$ in the left child $I_L$ using $O(\log t)$ bits. Bob compares this value with his count of $H$ in $I_L$. Since $H\subseteq F$, if the difference is positive recurse left; otherwise, under the invariant that $F\setminus H$ meets $I$, recurse right. At a singleton interval output its coordinate. Future behavior depends only on the current interval after the comparison, so all histories reaching the same interval merge. There are $O(N)$ intervals and at most $t+1$ count values; binary transmission adds the stated logarithmic factor.
\end{proof}

Hence the most natural assignment-witness KW restriction is \RET as a circuit-size lower-bound target. A successful DAG-like lower bound must exploit the full pair $(F,G)$ structure without giving Bob the entire satisfied part of $F$.

\section{TESTS AND OUTCOME}
A numerical rank audit formed the assignment-block Gram matrix $G(a,b)=\binom{q-d_H(a,b)}3$ and its entrywise powers. For $q=3,\ldots,8$ and $r=1,2,3$, the observed ranks were respectively: $q=3:(8,8,8)$; $4:(15,16,16)$; $5:(26,32,32)$; $6:(42,64,64)$; $7:(64,127,128)$; $8:(93,247,256)$. Every value equals $\sum_{i\le\min(3r,q)}\binom qi$. This audit is \COMP; the Fourier proof above is symbolic.

\paragraph{Outcome.}
The final monotone-slice task is now an exact combinatorial model: lower-bound the union/intersection set-circuit complexity of size-$3q$ transversals of the assignment-block hypergraph. Full-cube $\mathrm{UNSAT}_{3,q}$ hardness does not automatically transfer to this layer. HD tensors give an exact hierarchy of assignment information but no gate-stable superpolynomial potential. The easy masked KW protocol rules out the most immediate assignment-witness communication reduction. The surviving high-value targets are (i) a mass-production-resistant set-circuit invariant for the transversal family, or (ii) a lower bound for the full rectangle-DAG KW relation.

\paragraph{Next target.}
Study restrictions of the set circuit that preserve the transversal geometry while defeating Hiltgen--Paterson mass production. In particular, search for a family of assignment blocks whose restriction turns every low-gate set circuit into a bounded-union-spectrum object, allowing the moment--Bonferroni $\Xi$ bound to re-enter at the gate level. In parallel, audit whether Dunne's SAT central-slice projection has any set-circuit-preserving analogue for clause-presence UNSAT; ordinary NP-completeness alone is insufficient.

Root \texttt{main.tex} remains truncated and was not modified.

\begin{thebibliography}{7}
\bibitem{Wegener85} I. Wegener, \emph{On the Complexity of Slice Functions}, Theor. Comput. Sci. 38:55--68, 1985.
\bibitem{Wegener87} I. Wegener, \emph{The Complexity of Boolean Functions}, Wiley--Teubner, 1987, Sections 6.13--6.15.
\bibitem{HP95} A. P. Hiltgen and M. S. Paterson, \emph{PI$_k$ Mass Production and an Optimal Circuit for the Ne\v{c}iporuk Slice}, Computational Complexity 5:132--154, 1995.
\bibitem{Dunne86} P. E. Dunne, \emph{The Complexity of Central Slice Functions}, Theor. Comput. Sci. 44:247--257, 1986.
\bibitem{HP17} P. Hrube\v{s} and P. Pudl\'ak, \emph{Random Formulas, Monotone Circuits, and Interpolation}, FOCS 2017.
\bibitem{Sokolov17} D. Sokolov, \emph{Dag-like Communication and Its Applications}, CSR 2017; ECCC TR16-202.
\bibitem{Valiant86} L. G. Valiant, \emph{Negation Is Powerless for Boolean Slice Functions}, SIAM J. Comput. 15(2):531--535, 1986.
\end{thebibliography}
\end{document}
% FINAL_VERIFICATION_SENTINEL::MAIN::SLICESETTRANS1::2026-10-16_0915::P_VS_NP_OPEN